% GENERATED FILE. Edit canonical lessons, then run npm run generate:books.
% canonical-source-hash: fnv1a64:323d7c66a90ce4e0
% canonical-lessons: TA-C529-putir, TA-C529-vitukatai, TA-C529-tayam, TA-C529-pakatai, TA-C529-pallankuli, TA-R529-review-522-529

\chapter{Puzzles and Old Games}
\label{ch:ta-puzzles-and-old-games}
\hlchaptermodality{529}

\begin{chapteropening}
\textbf{By the end of this chapter:} \emph{I can name puzzles and old board games.}

Complete the last lesson of chapter 529: I can name puzzles and old board games.
\end{chapteropening}

\section[\texorpdfstring{\ta{புதிர்} (putir)}{புதிர் (putir)}]{\ta{புதிர்} (putir) — a puzzle}

\label{lesson:TA-C529-putir}



\begin{quote}
Before the new one: say the Tamil for to print and publish, then the Tamil for to compile.
\end{quote}

\subsection*{You'll want to know: \ta{புதிர்}}
\textbf{\ta{புதிர்}} — \emph{putir} — \textquotedblleft{}a puzzle\textquotedblright{}.

\begin{quote}
\ta{தாத்தா} \ta{ஒவ்வொரு} \ta{நாளும்} \ta{நாளிதழில்} \ta{வரும்} \ta{புதிருக்கு} \ta{விடை} \ta{கண்டுபிடிக்கிறார்}.
\end{quote}

\begin{quote}
\emph{tāttā ovvoru nāḷum nāḷitaḻil varum putirukku viṭai kaṇṭupiṭikkiṟār.}
\end{quote}

\begin{quote}
Every day grandfather solves the puzzle in the newspaper.
\end{quote}

To set someone a puzzle is \ta{புதிர்} \ta{போடு} (\emph{putir pōṭu}); to solve one is to find its answer: \ta{புதிருக்கு} \ta{விடை} \ta{கண்டுபிடி}.

\begin{cousinweb}
\textbf{\ta{புதிர்}} (\emph{putir}) is native Dravidian: DEDR 4163 gives \ta{பிதிர்} and \ta{புதிர்} (\emph{putir}), \textquotedblleft{}a riddle, a puzzle\textquotedblright{}.
\end{cousinweb}

\begin{grammarlens}[title={What you've built}]
One more word for describing what you lived and felt.
\end{grammarlens}

\subsection*{Your turn}
\begin{itemize}
\raggedright
  \item \emph{Say it:} \emph{putir}
  \item \emph{Say it:} \emph{putir}, once more
  \item \emph{Say it:} \emph{tāttā ovvoru nāḷum nāḷitaḻil varum putirukku viṭai kaṇṭupiṭikkiṟār.} — Every day grandfather solves the puzzle in the newspaper.
  \item \emph{Recall:} say the Tamil for to print and publish, then the Tamil for to compile, then say \emph{putir} again
\end{itemize}

\begin{tcolorbox}[breakable,colback=teal!4,colframe=teal!35!black,title={Before you move on}]
What does \ta{புதிர்} mean? (\textquotedblleft{}A puzzle\textquotedblright{}.) Say it once more.
\end{tcolorbox}
\section[\texorpdfstring{\ta{விடுகதை} (viṭukatai)}{விடுகதை (viṭukatai)}]{\ta{விடுகதை} (viṭukatai) — a riddle}

\label{lesson:TA-C529-vitukatai}



\begin{quote}
Before the new one: say the Tamil for to compile, then the Tamil for a puzzle.
\end{quote}

\subsection*{You'll want to know: \ta{விடுகதை}}
\textbf{\ta{விடுகதை}} — \emph{viṭukatai} — \textquotedblleft{}a riddle\textquotedblright{}.

\begin{quote}
\ta{பாட்டி} \ta{எங்களுக்கு} \ta{ஒரு} \ta{விடுகதை} \ta{சொன்னாள்}.
\end{quote}

\begin{quote}
\emph{pāṭṭi eṅkaḷukku oru viṭukatai coṉṉāḷ.}
\end{quote}

\begin{quote}
Grandmother told us a riddle.
\end{quote}

A spoken riddle, given as a little story; the listener must find the answer.

\begin{cousinweb}
\textbf{\ta{விடுகதை}} (\emph{viṭukatai}) ends in \ta{கதை} (\emph{katai}), \textquotedblleft{}a story\textquotedblright{}, from Sanskrit \emph{kathā} (CDIAL 2704); how the first part, \ta{விடு}, joins it is not settled.
\end{cousinweb}

\begin{grammarlens}[title={What you've built}]
One more word for describing what you lived and felt.
\end{grammarlens}

\subsection*{Your turn}
\begin{itemize}
\raggedright
  \item \emph{Say it:} \emph{viṭukatai}
  \item \emph{Say it:} \emph{viṭukatai}, once more
  \item \emph{Say it:} \emph{pāṭṭi eṅkaḷukku oru viṭukatai coṉṉāḷ.} — Grandmother told us a riddle.
  \item \emph{Recall:} say the Tamil for to compile, then the Tamil for a puzzle, then say \emph{viṭukatai} again
\end{itemize}

\begin{tcolorbox}[breakable,colback=teal!4,colframe=teal!35!black,title={Before you move on}]
What does \ta{விடுகதை} mean? (\textquotedblleft{}A riddle\textquotedblright{}.) Say it once more.
\end{tcolorbox}
\section[\texorpdfstring{\ta{தாயம்} (tāyam)}{தாயம் (tāyam)}]{\ta{தாயம்} (tāyam) — a dice board game}

\label{lesson:TA-C529-tayam}



\begin{quote}
Before the new one: say the Tamil for a puzzle, then the Tamil for a riddle.
\end{quote}

\subsection*{You'll want to know: \ta{தாயம்}}
\textbf{\ta{தாயம்}} — \emph{tāyam} — \textquotedblleft{}a dice board game\textquotedblright{}.

\begin{quote}
\ta{விடுமுறையில்} \ta{நாங்கள்} \ta{தாயம்} \ta{விளையாடினோம்}.
\end{quote}

\begin{quote}
\emph{viṭumuṟaiyil nāṅkaḷ tāyam viḷaiyāṭiṉōm.}
\end{quote}

\begin{quote}
In the holidays we played the dice game.
\end{quote}

A board game played with dice; the game takes \ta{விளையாடு} (\emph{viḷaiyāṭu}) like any other.

\begin{cousinweb}
The origin of \textbf{\ta{தாயம்}} (\emph{tāyam}) is not settled; one suggestion is Sanskrit \emph{dāya}, \textquotedblleft{}a share, a portion\textquotedblright{} (CDIAL 6288).
\end{cousinweb}

\begin{grammarlens}[title={What you've built}]
One more word for describing what you lived and felt.
\end{grammarlens}

\subsection*{Your turn}
\begin{itemize}
\raggedright
  \item \emph{Say it:} \emph{tāyam}
  \item \emph{Say it:} \emph{tāyam}, once more
  \item \emph{Say it:} \emph{viṭumuṟaiyil nāṅkaḷ tāyam viḷaiyāṭiṉōm.} — In the holidays we played the dice game.
  \item \emph{Recall:} say the Tamil for a puzzle, then the Tamil for a riddle, then say \emph{tāyam} again
\end{itemize}

\begin{tcolorbox}[breakable,colback=teal!4,colframe=teal!35!black,title={Before you move on}]
What does \ta{தாயம்} mean? (\textquotedblleft{}A dice board game\textquotedblright{}.) Say it once more.
\end{tcolorbox}
\section[\texorpdfstring{\ta{பகடை} (pakaṭai)}{பகடை (pakaṭai)}]{\ta{பகடை} (pakaṭai) — a die; dice}

\label{lesson:TA-C529-pakatai}



\begin{quote}
Before the new one: say the Tamil for a riddle, then the Tamil for a dice board game.
\end{quote}

\subsection*{You'll want to know: \ta{பகடை}}
\textbf{\ta{பகடை}} — \emph{pakaṭai} — \textquotedblleft{}a die; dice\textquotedblright{}.

\begin{quote}
\ta{அவன்} \ta{பகடையை} \ta{உருட்டினான்}.
\end{quote}

\begin{quote}
\emph{avaṉ pakaṭaiyai uruṭṭiṉāṉ.}
\end{quote}

\begin{quote}
He rolled the dice.
\end{quote}

One rolls dice with \ta{உருட்டு} (\emph{uruṭṭu}).

\begin{cousinweb}
The origin of \textbf{\ta{பகடை}} (\emph{pakaṭai}) is not settled: no dictionary we can check gives its source.
\end{cousinweb}

\begin{grammarlens}[title={What you've built}]
One more word for describing what you lived and felt.
\end{grammarlens}

\subsection*{Your turn}
\begin{itemize}
\raggedright
  \item \emph{Say it:} \emph{pakaṭai}
  \item \emph{Say it:} \emph{pakaṭai}, once more
  \item \emph{Say it:} \emph{avaṉ pakaṭaiyai uruṭṭiṉāṉ.} — He rolled the dice.
  \item \emph{Recall:} say the Tamil for a riddle, then the Tamil for a dice board game, then say \emph{pakaṭai} again
\end{itemize}

\begin{tcolorbox}[breakable,colback=teal!4,colframe=teal!35!black,title={Before you move on}]
What does \ta{பகடை} mean? (\textquotedblleft{}A die; dice\textquotedblright{}.) Say it once more.
\end{tcolorbox}
\section[\texorpdfstring{\ta{பல்லாங்குழி} (pallāṅkuḻi)}{பல்லாங்குழி (pallāṅkuḻi)}]{\ta{பல்லாங்குழி} (pallāṅkuḻi) — a seed-and-pit board game}

\label{lesson:TA-C529-pallankuli}



\begin{quote}
Before the new one: say the Tamil for a dice board game, then the Tamil for a die.
\end{quote}

\subsection*{You'll want to know: \ta{பல்லாங்குழி}}
\textbf{\ta{பல்லாங்குழி}} — \emph{pallāṅkuḻi} — \textquotedblleft{}a seed-and-pit board game\textquotedblright{}.

\begin{quote}
\ta{பாட்டியும்} \ta{பேத்தியும்} \ta{பல்லாங்குழி} \ta{விளையாடுகிறார்கள்}.
\end{quote}

\begin{quote}
\emph{pāṭṭiyum pēttiyum pallāṅkuḻi viḷaiyāṭukiṟārkaḷ.}
\end{quote}

\begin{quote}
Grandmother and granddaughter are playing the pit game.
\end{quote}

Two players move seeds or shells through rows of small pits in a wooden board.

\begin{cousinweb}
\textbf{\ta{பல்லாங்குழி}} (\emph{pallāṅkuḻi}) ends in the native \ta{குழி} (\emph{kuḻi}), \textquotedblleft{}a pit\textquotedblright{}, for the hollows of the board; the origin of the first part is not settled.
\end{cousinweb}

\begin{grammarlens}[title={What you've built}]
One more word for describing what you lived and felt.
\end{grammarlens}

\subsection*{Your turn}
\begin{itemize}
\raggedright
  \item \emph{Say it:} \emph{pallāṅkuḻi}
  \item \emph{Say it:} \emph{pallāṅkuḻi}, once more
  \item \emph{Say it:} \emph{pāṭṭiyum pēttiyum pallāṅkuḻi viḷaiyāṭukiṟārkaḷ.} — Grandmother and granddaughter are playing the pit game.
  \item \emph{Recall:} say the Tamil for a dice board game, then the Tamil for a die, then say \emph{pallāṅkuḻi} again
\end{itemize}

\begin{tcolorbox}[breakable,colback=teal!4,colframe=teal!35!black,title={Before you move on}]
What does \ta{பல்லாங்குழி} mean? (\textquotedblleft{}A seed-and-pit board game\textquotedblright{}.) Say it once more.
\end{tcolorbox}
\section[(dialogue)]{Looking back: from An Old Village House to Puzzles and Old Games}

\label{lesson:TA-R529-review-522-529}



\begin{quote}
Say the Tamil for a courtyard and a raised veranda seat.
\end{quote}

\subsection*{Your turn}
Say these aloud:

\begin{itemize}
\raggedright
  \item a courtyard, a raised veranda seat, a veranda, a pillar, a gutter
  \item to screw, to hang up, a blockage, a drain, to filter
  \item a sieve, a mortar, a pestle, to pour out, to purify
  \item cooking gas, firewood, charcoal, to steam, to be cooked
  \item to temper spices, to sauté, to stir a pot, to grate, to sprinkle
  \item fun, entertainment, to chat, a club, to support
  \item a writer, a novel, a poem, to print and publish, to compile
  \item a puzzle, a riddle, a dice board game, a die, a seed-and-pit board game
\end{itemize}

\begin{tcolorbox}[breakable,colback=teal!4,colframe=teal!35!black,title={Before you move on}]
A courtyard? (\textbf{\ta{முற்றம்}}, \emph{muṟṟam}.) A raised veranda seat? (\textbf{\ta{திண்ணை}}, \emph{tiṇṇai}.) A veranda? (\textbf{\ta{தாழ்வாரம்}}, \emph{tāḻvāram}.) A pillar? (\textbf{\ta{தூண்}}, \emph{tūṇ}.) A gutter? (\textbf{\ta{சாக்கடை}}, \emph{cākkaṭai}.) To screw? (\textbf{\ta{திருகு}}, \emph{tiruku}.) To hang up? (\textbf{\ta{மாட்டு}}, \emph{māṭṭu}.) A blockage? (\textbf{\ta{அடைப்பு}}, \emph{aṭaippu}.) A drain? (\textbf{\ta{வடிகால்}}, \emph{vaṭikāl}.) To filter? (\textbf{\ta{வடிகட்டு}}, \emph{vaṭikaṭṭu}.) A sieve? (\textbf{\ta{சல்லடை}}, \emph{callaṭai}.) A mortar? (\textbf{\ta{உரல்}}, \emph{ural}.) A pestle? (\textbf{\ta{உலக்கை}}, \emph{ulakkai}.) To pour out? (\textbf{\ta{கொட்டு}}, \emph{koṭṭu}.) To purify? (\textbf{\ta{சுத்திகரி}}, \emph{cuttikari}.) Cooking gas? (\textbf{\ta{எரிவாயு}}, \emph{erivāyu}.) Firewood? (\textbf{\ta{விறகு}}, \emph{viṟaku}.) Charcoal? (\textbf{\ta{கரி}}, \emph{kari}.) To steam? (\textbf{\ta{அவி}}, \emph{avi}.) To be cooked? (\textbf{\ta{வேகு}}, \emph{vēku}.) To temper spices? (\textbf{\ta{தாளி}}, \emph{tāḷi}.) To sauté? (\textbf{\ta{வதக்கு}}, \emph{vatakku}.) To stir a pot? (\textbf{\ta{கிளறு}}, \emph{kiḷaṟu}.) To grate? (\textbf{\ta{துருவு}}, \emph{turuvu}.) To sprinkle? (\textbf{\ta{தூவு}}, \emph{tūvu}.) Fun? (\textbf{\ta{வேடிக்கை}}, \emph{vēṭikkai}.) Entertainment? (\textbf{\ta{கேளிக்கை}}, \emph{kēḷikkai}.) To chat? (\textbf{\ta{அரட்டை} \ta{அடி}}, \emph{araṭṭai aṭi}.) A club? (\textbf{\ta{மன்றம்}}, \emph{maṉṟam}.) To support? (\textbf{\ta{ஆதரி}}, \emph{ātari}.) A writer? (\textbf{\ta{எழுத்தாளர்}}, \emph{eḻuttāḷar}.) A novel? (\textbf{\ta{நாவல்}}, \emph{nāval}.) A poem? (\textbf{\ta{கவிதை}}, \emph{kavitai}.) To print and publish? (\textbf{\ta{பதிப்பி}}, \emph{patippi}.) To compile? (\textbf{\ta{தொகு}}, \emph{toku}.) A puzzle? (\textbf{\ta{புதிர்}}, \emph{putir}.) A riddle? (\textbf{\ta{விடுகதை}}, \emph{viṭukatai}.) A dice board game? (\textbf{\ta{தாயம்}}, \emph{tāyam}.) A die? (\textbf{\ta{பகடை}}, \emph{pakaṭai}.) A seed-and-pit board game? (\textbf{\ta{பல்லாங்குழி}}, \emph{pallāṅkuḻi}.)
\end{tcolorbox}
